\section{组内课题范式：医学影像 AI 问题拆解}

本章抽象自组内「医学影像辅助检测」类课题梳理经验，重点示范\textbf{如何把模糊课题拆成可研究、可验证的子问题}，不绑定具体项目名称或未公开数据。

\subsection{从临床诉求到系统定位}
典型临床痛点包括：阅片疲劳导致漏诊、微小/隐匿病灶难发现、报告标准不统一。\\
对应系统目标常包括：提高检出、缩短报告时间、统一标准、控制假阳性。\\
入门时先分清能力层级，例如：
\begin{enumerate}
  \item \textbf{提质增效}：阅片前自动检测与标记（CADe）——多数入门课题先落在此层；
  \item \textbf{智能预警}：高风险病灶优先级提醒；
  \item \textbf{高级辅助}：诊断建议、结构化报告、复杂决策支持。
\end{enumerate}
写开题/汇报时明确边界：辅助检出 $\neq$ 替代诊断。

\subsection{技术本质常见画像}
医学影像检测问题常同时具备：
\begin{itemize}
  \item \textbf{小目标}：病灶尺度小、信号弱；
  \item \textbf{异常检测}：低对比、位置隐匿；
  \item \textbf{极度类别不平衡}：正常体素远多于病灶。
\end{itemize}
不同病种在模态（CT/MRI/X-ray）、维度（2D/3D）与外观上差异大，可能需要多任务学习或病种专用检测器——选题时先做数据与任务摸底。

\subsection{推荐的子问题拆解框架}
可将 pipeline 正交拆成若干可独立调研的子问题，例如：
\begin{enumerate}
  \item \textbf{弱信号保留}：下采样中如何保住微小病灶特征？多尺度如何融合？
  \item \textbf{召回--假阳性权衡}：高召回下如何把 FPR 压到临床可接受区间？
  \item \textbf{多模态语义}：病史/报告文本与影像如何对齐？结构化输出如何保证事实准确？
  \item \textbf{工程部署}：精度与时延/边缘算力如何权衡？如何接入现有工作流？
  \item \textbf{数据与标注}：标注稀缺、不一致、不平衡如何缓解（弱监督/半监督等）？
  \item \textbf{评价体系}：如何定义并衡量「诊断质量提升」，而不只是刷单一指标？
\end{enumerate}
\tipbox{先把「问题」与「具体网络结构方案」分开写。级联网络、某损失等是手段，不应在问题尚未清晰时就锁死方案。}

\subsection{入门四周可执行的下一步}
\begin{enumerate}
  \item 针对你被分配的 1--2 个子问题，各检索若干 SOTA 与开源基准；
  \item 数据摸底：公开/可申请数据、标注规模与许可；
  \item 与导师确认首个切入点与成功标准（指标与临床约束）；
  \item 用第~3~章模板产出文献 list，并做一次 15--20 分钟汇报。
\end{enumerate}

\subsection{数据集检索入口}
\begin{itemize}
  \item Nature Scientific Data：\url{https://www.nature.com/sdata/}
  \item Grand Challenge（医学影像竞赛数据）：\url{https://grand-challenge.org/}
  \item Kaggle Datasets：\url{https://www.kaggle.com/datasets}
  \item Papers with Code 数据集页；以及从顶会论文中挖作者公开数据。
\end{itemize}
使用前务必阅读许可协议、获取门槛与标注说明。数据集整理字段模板见 \texttt{templates/dataset\_sheet.csv}。
